% ============================================================
% 讲义 — 人教版高中必修第一册 · 第2章 一元二次函数、方程和不等式
% 编译命令：xelatex -interaction=nonstopmode 高一必修1_第2章_一元二次函数方程和不等式.tex
% ============================================================
\documentclass[12pt, a4paper]{ctexart}
\usepackage[margin=2.2cm]{geometry}
\usepackage{amsmath, amssymb, amsthm}
\usepackage{enumitem}
\usepackage{tcolorbox}
\usepackage{booktabs}
\usepackage{array}
\usepackage{fancyhdr}
\usepackage{xcolor}
\usepackage{tikz}
\usepackage{hyperref}
\usepackage{graphicx}
\usepackage{multirow}
\usepackage{pgfplots}
\pgfplotsset{compat=1.18}

% 颜色定义
\definecolor{mainblue}{RGB}{41, 98, 168}
\definecolor{lightblue}{RGB}{230, 240, 255}
\definecolor{tipgreen}{RGB}{34, 139, 87}
\definecolor{lightgreen}{RGB}{235, 250, 240}
\definecolor{warnred}{RGB}{200, 50, 50}
\definecolor{lightred}{RGB}{255, 240, 240}
\definecolor{examorange}{RGB}{210, 120, 20}
\definecolor{lightorange}{RGB}{255, 248, 235}
\definecolor{apppurple}{RGB}{120, 60, 170}
\definecolor{lightpurple}{RGB}{245, 238, 255}

% 页眉页脚
\pagestyle{fancy}
\fancyhf{}
\fancyhead[L]{\small\textcolor{mainblue}{\textbf{人教版高中必修第一册 · 第2章 一元二次函数、方程和不等式}}}
\fancyhead[R]{\small\textcolor{mainblue}{高一数学讲义}}
\fancyfoot[C]{\small —\ \thepage\ —}
\renewcommand{\headrulewidth}{0.6pt}
\renewcommand{\headrule}{\hbox to\headwidth{\color{mainblue}\leaders\hrule height \headrulewidth\hfill}}
\renewcommand{\footrulewidth}{0pt}

% tcolorbox 样式
\newtcolorbox{knowledgebox}[1][]{%
  colback=lightblue, colframe=mainblue,
  fonttitle=\bfseries\large, title=#1,
  arc=3mm, boxrule=0.8pt,
  left=8pt, right=8pt, top=6pt, bottom=6pt
}

\newtcolorbox{tipbox}[1][]{%
  colback=lightgreen, colframe=tipgreen,
  fonttitle=\bfseries, title=#1,
  arc=2mm, boxrule=0.6pt,
  left=6pt, right=6pt, top=4pt, bottom=4pt
}

\newtcolorbox{warnbox}[1][]{%
  colback=lightred, colframe=warnred,
  fonttitle=\bfseries, title=#1,
  arc=2mm, boxrule=0.6pt,
  left=6pt, right=6pt, top=4pt, bottom=4pt
}

\newtcolorbox{examplebox}[1][]{%
  colback=lightorange, colframe=examorange,
  fonttitle=\bfseries, title=#1,
  arc=3mm, boxrule=0.8pt,
  left=8pt, right=8pt, top=6pt, bottom=6pt
}

\newtcolorbox{appbox}[1][]{%
  colback=lightpurple, colframe=apppurple,
  fonttitle=\bfseries, title=#1,
  arc=3mm, boxrule=0.8pt,
  left=8pt, right=8pt, top=6pt, bottom=6pt
}

\setCJKmainfont{Hiragino Sans GB}
\setCJKsansfont{Heiti SC}

\begin{document}

% ============ 标题页 ============
\begin{center}
{\Huge\bfseries\textcolor{mainblue}{第2章\quad 一元二次函数、方程和不等式}}\\[12pt]
{\Large\textcolor{gray}{高一数学讲义}}\\[8pt]
\rule{0.6\textwidth}{1pt}\\[8pt]
{\normalsize 适用：人教版高中必修第一册\quad |\quad 课时：10课时}
\end{center}
\vspace{10pt}

% =====================================================
% 第一部分：知识点讲解
% =====================================================
\section*{\textcolor{mainblue}{一、知识点讲解}}

% ──────── 2.1 等式性质与不等式性质 ────────
\subsection*{2.1\quad 等式性质与不等式性质}

\begin{knowledgebox}[等式的基本性质]
\begin{enumerate}[leftmargin=2em]
  \item 若 $a = b$，则 $b = a$（对称性）。
  \item 若 $a = b$，$b = c$，则 $a = c$（传递性）。
  \item 若 $a = b$，则 $a + c = b + c$。
  \item 若 $a = b$，则 $ac = bc$。
  \item 若 $a = b$，$c = d$，则 $a + c = b + d$，$ac = bd$。
\end{enumerate}
\end{knowledgebox}

\begin{knowledgebox}[不等式的基本性质]
\begin{enumerate}[leftmargin=2em]
  \item \textbf{对称性}：$a > b \Leftrightarrow b < a$。
  \item \textbf{传递性}：$a > b$ 且 $b > c \Rightarrow a > c$。
  \item \textbf{加法单调性}：$a > b \Rightarrow a + c > b + c$。\\
    推论：$a + b > c \Leftrightarrow a > c - b$（不等式可以移项）。
  \item \textbf{乘法性质}：
    \begin{itemize}
      \item $a > b,\ c > 0 \Rightarrow ac > bc$（同乘正数不变号）。
      \item $a > b,\ c < 0 \Rightarrow ac < bc$（同乘负数要变号）。
    \end{itemize}
  \item \textbf{同向可加性}：$a > b$ 且 $c > d \Rightarrow a + c > b + d$。\\
    注意：\textbf{同向不等式不能相减}！（如 $3 > 1$，$5 > 2$，但 $3 - 5 > 1 - 2$ 即 $-2 > -1$ 不成立。）
  \item \textbf{同向同正可乘}：$a > b > 0$ 且 $c > d > 0 \Rightarrow ac > bd$。
  \item \textbf{乘方}：$a > b > 0 \Rightarrow a^n > b^n$（$n \in \mathbb{N}^*$）。
  \item \textbf{开方}：$a > b > 0 \Rightarrow \sqrt[n]{a} > \sqrt[n]{b}$（$n \in \mathbb{N}^*,\ n \geq 2$）。
\end{enumerate}
\end{knowledgebox}

\begin{warnbox}[易错提醒：不等式运算]
\begin{itemize}[leftmargin=2em]
  \item 不等式\textbf{两边同乘负数要变号}！例如 $-2x > 4$，两边除以 $-2$ 得 $x < -2$。
  \item 不等式\textbf{不能随便相除}。$a > b > 0$ 且 $c > d > 0$ 不能推出 $\dfrac{a}{c} > \dfrac{b}{d}$。例如 $6 > 2$，$4 > 1$，但 $\dfrac{6}{4} < \dfrac{2}{1}$。
  \item 不等式两边平方需谨慎：$a > b$ 不能推出 $a^2 > b^2$（除非 $a > b > 0$ 或 $a > b$ 且 $|a| > |b|$）。
\end{itemize}
\end{warnbox}

% ──────── 2.2 基本不等式 ────────
\subsection*{2.2\quad 基本不等式}

\begin{knowledgebox}[基本不等式（均值不等式）]
若 $a > 0$，$b > 0$，则：
\[
  \frac{a + b}{2} \geq \sqrt{ab}
\]
即
\[
  a + b \geq 2\sqrt{ab}
\]
等号成立当且仅当 $a = b$。
\begin{itemize}[leftmargin=2em]
  \item $\dfrac{a+b}{2}$ 叫做 $a, b$ 的\textbf{算术平均数}。
  \item $\sqrt{ab}$ 叫做 $a, b$ 的\textbf{几何平均数}。
  \item 基本不等式可表述为：\textbf{两个正数的算术平均数不小于几何平均数}。
\end{itemize}
\end{knowledgebox}

\begin{knowledgebox}[基本不等式的证明]
\textbf{证明}：因为 $a > 0$，$b > 0$，所以：
\[
  a + b - 2\sqrt{ab} = (\sqrt{a})^2 - 2\sqrt{a}\sqrt{b} + (\sqrt{b})^2 = (\sqrt{a} - \sqrt{b})^2 \geq 0
\]
因此 $a + b \geq 2\sqrt{ab}$。等号成立当且仅当 $\sqrt{a} = \sqrt{b}$，即 $a = b$。\qed
\end{knowledgebox}

\begin{knowledgebox}[基本不等式的常见变形与推论]
\begin{enumerate}[leftmargin=2em]
  \item $ab \leq \left(\dfrac{a+b}{2}\right)^2$（已知 $a + b$ 求 $ab$ 的最大值）。
  \item 若 $a > 0$，则 $a + \dfrac{1}{a} \geq 2$，等号在 $a = 1$ 时成立。
  \item 若 $ab > 0$，则 $\dfrac{a}{b} + \dfrac{b}{a} \geq 2$，等号在 $a = b$ 时成立。
  \item 若 $a > 0,\ b > 0$，则 $a^2 + b^2 \geq 2ab$。
\end{enumerate}
\end{knowledgebox}

\begin{warnbox}[使用基本不等式的三个条件（一正二定三相等）]
使用 $a + b \geq 2\sqrt{ab}$ 求最值时，必须同时满足：
\begin{enumerate}[leftmargin=2em]
  \item \textbf{一正}：$a > 0$，$b > 0$（各项均为正数）。
  \item \textbf{二定}：$a + b$（或 $ab$）为定值。
  \item \textbf{三相等}：等号能够取到（即 $a = b$ 的条件可以满足）。
\end{enumerate}
\textbf{三个条件缺一不可！}
\end{warnbox}

\begin{tipbox}[基本不等式的应用技巧]
\textbf{求和的最小值}：若 $ab = $ 常数，则 $a + b \geq 2\sqrt{ab}$，当 $a = b$ 时取最小值。\\
\textbf{求积的最大值}：若 $a + b = $ 常数，则 $ab \leq \left(\dfrac{a+b}{2}\right)^2$，当 $a = b$ 时取最大值。\\
\textbf{拆项凑定值}：如求 $x + \dfrac{4}{x}$（$x > 0$）的最小值，直接用 $x + \dfrac{4}{x} \geq 2\sqrt{x \cdot \dfrac{4}{x}} = 4$，等号在 $x = 2$ 时成立。
\end{tipbox}

% ──────── 2.3 一元二次函数的图象与性质 ────────
\subsection*{2.3\quad 一元二次函数的图象与性质}

\begin{knowledgebox}[一元二次函数]
一元二次函数的一般形式为 $f(x) = ax^2 + bx + c$（$a \neq 0$）。

\textbf{顶点式}：$f(x) = a\left(x + \dfrac{b}{2a}\right)^2 + \dfrac{4ac - b^2}{4a}$，顶点坐标为 $\left(-\dfrac{b}{2a},\ \dfrac{4ac - b^2}{4a}\right)$。
\end{knowledgebox}

\begin{knowledgebox}[一元二次函数的性质]
设 $f(x) = ax^2 + bx + c$（$a \neq 0$），令 $\Delta = b^2 - 4ac$。
\begin{enumerate}[leftmargin=2em]
  \item \textbf{开口方向}：$a > 0$ 开口向上，$a < 0$ 开口向下。
  \item \textbf{对称轴}：$x = -\dfrac{b}{2a}$。
  \item \textbf{顶点}：$\left(-\dfrac{b}{2a},\ \dfrac{4ac - b^2}{4a}\right) = \left(-\dfrac{b}{2a},\ -\dfrac{\Delta}{4a}\right)$。
  \item \textbf{与 $x$ 轴的交点}（即 $f(x) = 0$ 的解）：
    \begin{itemize}
      \item $\Delta > 0$：两个不同的实数根 $x_{1,2} = \dfrac{-b \pm \sqrt{\Delta}}{2a}$，图象与 $x$ 轴有两个交点。
      \item $\Delta = 0$：一个实数根 $x = -\dfrac{b}{2a}$，图象与 $x$ 轴相切。
      \item $\Delta < 0$：无实数根，图象与 $x$ 轴无交点。
    \end{itemize}
  \item \textbf{与 $y$ 轴的交点}：$(0, c)$。
  \item \textbf{单调性}（$a > 0$ 时）：在 $\left(-\infty,\ -\dfrac{b}{2a}\right)$ 上单调递减，在 $\left(-\dfrac{b}{2a},\ +\infty\right)$ 上单调递增。\\
    $a < 0$ 时相反。
\end{enumerate}
\end{knowledgebox}

\begin{tipbox}[二次函数图象速记]
$a$ 的正负决定开口方向，$-\dfrac{b}{2a}$ 决定对称轴位置，$c$ 决定与 $y$ 轴交点。\\
画图时先找\textbf{开口方向}、\textbf{对称轴}、\textbf{顶点}、\textbf{与坐标轴的交点}，再用对称性画出另一半。
\end{tipbox}

% ──────── 2.4 一元二次方程的解法 ────────
\subsection*{2.4\quad 一元二次方程的解法}

\begin{knowledgebox}[一元二次方程 $ax^2 + bx + c = 0$ 的解法]
\begin{enumerate}[leftmargin=2em]
  \item \textbf{直接开平方法}：形如 $ax^2 + c = 0$ 时，$x^2 = -\dfrac{c}{a}$，$x = \pm\sqrt{-\dfrac{c}{a}}$。
  \item \textbf{配方法}：将 $ax^2 + bx + c = 0$ 配方为 $\left(x + \dfrac{b}{2a}\right)^2 = \dfrac{b^2 - 4ac}{4a^2}$，再开方求解。
  \item \textbf{公式法}：$x = \dfrac{-b \pm \sqrt{b^2 - 4ac}}{2a}$（$\Delta \geq 0$ 时有实数根）。
  \item \textbf{因式分解法}：若 $ax^2 + bx + c = a(x - x_1)(x - x_2)$，则 $x = x_1$ 或 $x = x_2$。
\end{enumerate}
\end{knowledgebox}

\begin{knowledgebox}[判别式 $\Delta = b^2 - 4ac$ 的作用]
\begin{center}
\renewcommand{\arraystretch}{1.3}
\begin{tabular}{c|c|c}
\toprule
$\Delta$ 的符号 & 根的情况 & 图象与 $x$ 轴的关系 \\
\midrule
$\Delta > 0$ & 两个不相等的实数根 & 两个交点 \\
$\Delta = 0$ & 两个相等的实数根（重根） & 一个交点（相切） \\
$\Delta < 0$ & 无实数根 & 无交点 \\
\bottomrule
\end{tabular}
\end{center}
\end{knowledgebox}

\begin{knowledgebox}[韦达定理（根与系数的关系）]
若 $ax^2 + bx + c = 0$（$a \neq 0$）的两个根为 $x_1, x_2$，则：
\[
  x_1 + x_2 = -\frac{b}{a}, \qquad x_1 \cdot x_2 = \frac{c}{a}
\]
\textbf{应用}：已知两根之和、两根之积，可以反过来构造方程：$x^2 - (x_1+x_2)x + x_1 x_2 = 0$。
\end{knowledgebox}

% ──────── 2.5 一元二次不等式的解法 ────────
\subsection*{2.5\quad 一元二次不等式的解法（核心）}

\begin{knowledgebox}[一元二次不等式]
一元二次不等式的标准形式：
\[
  ax^2 + bx + c > 0 \quad \text{或} \quad ax^2 + bx + c < 0 \quad (a \neq 0)
\]
\textbf{解题步骤}：
\begin{enumerate}[leftmargin=2em]
  \item 将不等式化为标准形式（$a > 0$，若 $a < 0$ 两边乘 $-1$ 并变号）。
  \item 计算判别式 $\Delta = b^2 - 4ac$。
  \item 根据 $\Delta$ 的符号，分三种情况讨论（见下表）。
\end{enumerate}
\end{knowledgebox}

\begin{knowledgebox}[$a > 0$ 时三种情况的解集]
设 $a > 0$，$\Delta = b^2 - 4ac$，$x_1 < x_2$ 为方程 $ax^2 + bx + c = 0$ 的两个根。

\begin{center}
\renewcommand{\arraystretch}{1.6}
\begin{tabular}{c|c|c}
\toprule
& $ax^2 + bx + c > 0$ 的解集 & $ax^2 + bx + c < 0$ 的解集 \\
\midrule
$\Delta > 0$（两根 $x_1 < x_2$） & $(-\infty,\ x_1) \cup (x_2,\ +\infty)$ & $(x_1,\ x_2)$ \\
$\Delta = 0$（重根 $x_0$） & $\{x \mid x \neq x_0\}$，即 $(-\infty, x_0) \cup (x_0, +\infty)$ & $\varnothing$ \\
$\Delta < 0$（无实根） & $\mathbb{R}$（全体实数） & $\varnothing$ \\
\bottomrule
\end{tabular}
\end{center}
\textbf{记忆口诀}："大于取两边，小于取中间"（$\Delta > 0$ 时）。
\end{knowledgebox}

\begin{warnbox}[易错提醒：解一元二次不等式]
\begin{itemize}[leftmargin=2em]
  \item \textbf{先化为 $a > 0$ 的形式}！若原题为 $-x^2 + 3x - 2 > 0$，先两边乘 $-1$ 变为 $x^2 - 3x + 2 < 0$（不等号变向！）。
  \item 注意\textbf{端点是否取等}：$>$ 和 $<$ 是严格不等号，端点不取；$\geq$ 和 $\leq$ 非严格，端点取。
  \item $\Delta = 0$ 时，$ax^2 + bx + c = a(x - x_0)^2 \geq 0$，只有 $x = x_0$ 时等于 $0$，其余均大于 $0$。
  \item 别忘了写成\textbf{集合或区间}的形式。
\end{itemize}
\end{warnbox}

\begin{tipbox}[一元二次不等式的图象理解法]
\textbf{画抛物线 $y = ax^2 + bx + c$ 的草图}，从图上直接读出 $y > 0$ 或 $y < 0$ 对应的 $x$ 的范围。\\
这是最直观、最不容易出错的方法！
\begin{itemize}[leftmargin=2em]
  \item $a > 0$，开口向上，$\Delta > 0$：抛物线与 $x$ 轴有两个交点 $x_1, x_2$。\\
    图象在 $x$ 轴上方（$y > 0$）的部分：$x < x_1$ 或 $x > x_2$。\\
    图象在 $x$ 轴下方（$y < 0$）的部分：$x_1 < x < x_2$。
  \item $a > 0$，$\Delta < 0$：抛物线完全在 $x$ 轴上方，$y > 0$ 对所有 $x$ 成立，$y < 0$ 无解。
\end{itemize}
\end{tipbox}

% ──────── 2.6 二次函数、方程、不等式的联系 ────────
\subsection*{2.6\quad 二次函数、方程、不等式的联系}

\begin{knowledgebox}[三者之间的关系]
\textbf{二次函数} $y = ax^2 + bx + c$ 是核心纽带：
\begin{itemize}[leftmargin=2em]
  \item \textbf{方程} $ax^2 + bx + c = 0$ 的根 $\Leftrightarrow$ 函数图象与 $x$ 轴交点的横坐标。
  \item \textbf{不等式} $ax^2 + bx + c > 0$ 的解集 $\Leftrightarrow$ 函数图象在 $x$ 轴上方的 $x$ 的范围。
  \item \textbf{不等式} $ax^2 + bx + c < 0$ 的解集 $\Leftrightarrow$ 函数图象在 $x$ 轴下方的 $x$ 的范围。
\end{itemize}
\textbf{关键}：画出二次函数的图象，方程和不等式的解一目了然！
\end{knowledgebox}

\begin{tipbox}[二次函数、方程、不等式的统一理解]
\begin{center}
\renewcommand{\arraystretch}{1.4}
\begin{tabular}{c|c|c|c}
\toprule
$\Delta$ & 方程 $= 0$ & 不等式 $> 0$ & 不等式 $< 0$ \\
\midrule
$\Delta > 0$ & $x_1,\ x_2$ & $(-\infty, x_1) \cup (x_2, +\infty)$ & $(x_1, x_2)$ \\
$\Delta = 0$ & $x_0$（重根） & $x \neq x_0$ & $\varnothing$ \\
$\Delta < 0$ & 无解 & $\mathbb{R}$ & $\varnothing$ \\
\bottomrule
\end{tabular}
\end{center}
以上均假设 $a > 0$。若 $a < 0$，$> 0$ 和 $< 0$ 的解集互换。
\end{tipbox}

% =====================================================
% 第二部分：例题精讲
% =====================================================
\section*{\textcolor{examorange}{二、例题精讲}}

\begin{examplebox}[【例1】解一元二次不等式]
解不等式 $x^2 - 5x + 6 < 0$。
\end{examplebox}

\textbf{解}：
\begin{enumerate}[leftmargin=2em]
  \item 已是标准形式 $a = 1 > 0$。
  \item 计算 $\Delta = 25 - 24 = 1 > 0$，方程 $x^2 - 5x + 6 = 0$ 有两个不等实根。
  \item 解方程：$(x-2)(x-3) = 0$，$x_1 = 2$，$x_2 = 3$。
  \item 因为 $a > 0$，$\Delta > 0$，"$< 0$"取中间：$x \in (2, 3)$。
\end{enumerate}
\textbf{答}：解集为 $\boxed{(2,\ 3)}$。

\begin{examplebox}[【例2】解含负系数的一元二次不等式]
解不等式 $-2x^2 + 3x + 5 \geq 0$。
\end{examplebox}

\textbf{解}：
\begin{enumerate}[leftmargin=2em]
  \item 两边乘 $-1$（变号）：$2x^2 - 3x - 5 \leq 0$。
  \item $\Delta = 9 + 40 = 49$，$\sqrt{\Delta} = 7$。
  \item $x = \dfrac{3 \pm 7}{4}$，$x_1 = \dfrac{3-7}{4} = -1$，$x_2 = \dfrac{3+7}{4} = \dfrac{5}{2}$。
  \item $a = 2 > 0$，"$\leq 0$"取中间（含端点）：$x \in \left[-1,\ \dfrac{5}{2}\right]$。
\end{enumerate}
\textbf{答}：解集为 $\boxed{\left[-1,\ \dfrac{5}{2}\right]}$。

\begin{examplebox}[【例3】$\Delta < 0$ 的情况]
解不等式 $x^2 + 2x + 5 > 0$。
\end{examplebox}

\textbf{解}：
\begin{enumerate}[leftmargin=2em]
  \item $a = 1 > 0$。
  \item $\Delta = 4 - 20 = -16 < 0$。
  \item 方程 $x^2 + 2x + 5 = 0$ 无实数根。
  \item 抛物线 $y = x^2 + 2x + 5$ 完全在 $x$ 轴上方，所以 $y > 0$ 对所有 $x$ 成立。
\end{enumerate}
\textbf{答}：解集为 $\boxed{\mathbb{R}}$。

\begin{tipbox}[方法总结]
$\Delta < 0$ 时：$a > 0$ 则 $ax^2 + bx + c > 0$ 的解集为 $\mathbb{R}$，$ax^2 + bx + c < 0$ 的解集为 $\varnothing$。\\
可以通过配方验证：$x^2 + 2x + 5 = (x+1)^2 + 4 \geq 4 > 0$。
\end{tipbox}

\begin{examplebox}[【例4】含参数的一元二次不等式]
已知不等式 $ax^2 + 2x + 1 > 0$ 的解集为 $\mathbb{R}$，求 $a$ 的取值范围。
\end{examplebox}

\textbf{解}：需要分情况讨论。
\begin{itemize}[leftmargin=2em]
  \item \textbf{若 $a = 0$}：不等式变为 $2x + 1 > 0$，解集为 $x > -\dfrac{1}{2}$，不是 $\mathbb{R}$，不满足。
  \item \textbf{若 $a \neq 0$}：要使 $ax^2 + 2x + 1 > 0$ 对所有 $x$ 成立，需要：
    \begin{enumerate}[leftmargin=2em]
      \item $a > 0$（开口向上）。
      \item $\Delta < 0$（图象完全在 $x$ 轴上方）。
    \end{enumerate}
    $\Delta = 4 - 4a < 0 \Rightarrow a > 1$。
\end{itemize}
\textbf{答}：$a$ 的取值范围为 $\boxed{(1,\ +\infty)}$。

\begin{examplebox}[【例5】利用基本不等式求最值]
\begin{enumerate}[label=(\arabic*)]
  \item 若 $x > 0$，求 $x + \dfrac{9}{x}$ 的最小值。
  \item 若 $x > 0,\ y > 0$，且 $x + y = 4$，求 $xy$ 的最大值。
\end{enumerate}
\end{examplebox}

\textbf{解}：
\begin{enumerate}[label=(\arabic*)]
  \item $x > 0$，$\dfrac{9}{x} > 0$，满足"一正"。\\
    $x + \dfrac{9}{x} \geq 2\sqrt{x \cdot \dfrac{9}{x}} = 2\sqrt{9} = 6$。\\
    等号条件：$x = \dfrac{9}{x}$，即 $x = 3$。满足"三相等"。\\
    最小值为 $\boxed{6}$（$x = 3$ 时取得）。
  \item $x > 0,\ y > 0$，$x + y = 4$ 为定值，满足"一正二定"。\\
    $xy \leq \left(\dfrac{x+y}{2}\right)^2 = \left(\dfrac{4}{2}\right)^2 = 4$。\\
    等号条件：$x = y = 2$。\\
    最大值为 $\boxed{4}$（$x = y = 2$ 时取得）。
\end{enumerate}

\begin{warnbox}[提醒：凑定值的技巧]
当直接使用基本不等式无法凑出定值时，需要\textbf{拆项}。例如：\\
若 $x > 1$，求 $x + \dfrac{4}{x - 1}$ 的最小值。\\
令 $t = x - 1 > 0$，则 $x + \dfrac{4}{x-1} = (t + 1) + \dfrac{4}{t} = t + \dfrac{4}{t} + 1 \geq 2\sqrt{4} + 1 = 5$。\\
等号在 $t = 2$ 即 $x = 3$ 时成立。最小值为 $5$。
\end{warnbox}

\begin{examplebox}[【例6】二次函数与方程根的关系]
已知二次函数 $f(x) = x^2 - 2x - 3$。
\begin{enumerate}[label=(\arabic*)]
  \item 求 $f(x) = 0$ 的根。
  \item 求 $f(x) > 0$ 和 $f(x) < 0$ 的解集。
  \item 画出 $f(x)$ 的图象。
\end{enumerate}
\end{examplebox}

\textbf{解}：
\begin{enumerate}[label=(\arabic*)]
  \item $x^2 - 2x - 3 = 0$，$(x-3)(x+1) = 0$，$x_1 = -1$，$x_2 = 3$。
  \item $a = 1 > 0$，$\Delta = 4 + 12 = 16 > 0$。
    \begin{itemize}
      \item $f(x) > 0$ 的解集：$(-\infty,\ -1) \cup (3,\ +\infty)$。
      \item $f(x) < 0$ 的解集：$(-1,\ 3)$。
    \end{itemize}
  \item 对称轴 $x = 1$，顶点 $(1, -4)$，与 $y$ 轴交点 $(0, -3)$，与 $x$ 轴交点 $(-1, 0)$ 和 $(3, 0)$。
\end{enumerate}

\begin{examplebox}[【例7】韦达定理的应用]
已知方程 $x^2 - 5x + 3 = 0$ 的两根为 $x_1, x_2$，不解方程，求 $x_1^2 + x_2^2$ 和 $\dfrac{1}{x_1} + \dfrac{1}{x_2}$ 的值。
\end{examplebox}

\textbf{解}：由韦达定理：$x_1 + x_2 = 5$，$x_1 x_2 = 3$。

\begin{enumerate}[leftmargin=2em]
  \item $x_1^2 + x_2^2 = (x_1 + x_2)^2 - 2x_1 x_2 = 25 - 6 = \boxed{19}$。
  \item $\dfrac{1}{x_1} + \dfrac{1}{x_2} = \dfrac{x_1 + x_2}{x_1 x_2} = \dfrac{5}{3}$。
\end{enumerate}

\begin{tipbox}[韦达定理常用公式]
\begin{itemize}[leftmargin=2em]
  \item $x_1^2 + x_2^2 = (x_1 + x_2)^2 - 2x_1 x_2$
  \item $|x_1 - x_2| = \sqrt{(x_1 + x_2)^2 - 4x_1 x_2} = \dfrac{\sqrt{\Delta}}{|a|}$
  \item $x_1^3 + x_2^3 = (x_1 + x_2)^3 - 3x_1 x_2(x_1 + x_2)$
\end{itemize}
\end{tipbox}

\begin{examplebox}[【例8】不等式的综合应用]
解关于 $x$ 的不等式 $x^2 - (a+1)x + a < 0$（$a$ 为常数）。
\end{examplebox}

\textbf{解}：先因式分解。$x^2 - (a+1)x + a = (x - 1)(x - a)$。

所以不等式为 $(x-1)(x-a) < 0$。

需要比较 $1$ 和 $a$ 的大小：
\begin{itemize}[leftmargin=2em]
  \item \textbf{若 $a > 1$}：两根 $x_1 = 1$，$x_2 = a$，$1 < a$。解集为 $(1,\ a)$。
  \item \textbf{若 $a = 1$}：$(x-1)^2 < 0$，无解。解集为 $\varnothing$。
  \item \textbf{若 $a < 1$}：两根 $x_1 = a$，$x_2 = 1$，$a < 1$。解集为 $(a,\ 1)$。
\end{itemize}
\textbf{答}：
\[
  \text{解集} = \begin{cases} (1,\ a) & \text{若 } a > 1 \\ \varnothing & \text{若 } a = 1 \\ (a,\ 1) & \text{若 } a < 1 \end{cases}
\]

% =====================================================
% 第三部分：应用题
% =====================================================
\section*{\textcolor{apppurple}{三、应用题}}

\begin{appbox}[【应用1】利润最大化问题]
某商店以每件 $40$ 元的价格购进一批商品。据市场调查，当售价为每件 $x$ 元（$x > 40$）时，每天能卖出 $(200 - 2x + 80) = (280 - 2x)$ 件。每件商品的利润为 $(x - 40)$ 元。
\begin{enumerate}[label=(\arabic*)]
  \item 写出每天总利润 $y$（元）关于售价 $x$（元）的函数。
  \item 当售价定为多少时，每天的总利润最大？最大利润是多少？
\end{enumerate}
\end{appbox}

\textbf{解}：
\begin{enumerate}[label=(\arabic*)]
  \item 每天总利润 = 每件利润 $\times$ 销量：
    \[
      y = (x - 40)(280 - 2x) = -2x^2 + 360x - 11200
    \]
    其中 $40 < x < 140$（售价不能超过 $140$ 元，否则销量为 $0$）。
  \item 这是一个开口向下的抛物线 $a = -2 < 0$，对称轴为 $x = -\dfrac{360}{2 \times (-2)} = 90$。\\
    当 $x = 90$ 时，$y_{\max} = -2 \times 90^2 + 360 \times 90 - 11200 = -16200 + 32400 - 11200 = \boxed{5000}$（元）。
\end{enumerate}
\textbf{答}：售价定为每件 $90$ 元时，每天利润最大，为 $5000$ 元。

\begin{appbox}[【应用2】围栏面积问题]
用总长为 $60$ 米的篱笆围成一面靠墙的矩形花园（墙的长度足够），设垂直于墙的一边长为 $x$ 米。
\begin{enumerate}[label=(\arabic*)]
  \item 写出花园面积 $S$ 关于 $x$ 的函数，并指出定义域。
  \item 当 $x$ 为何值时，花园面积最大？最大面积是多少？
\end{enumerate}
\end{appbox}

\textbf{解}：
\begin{enumerate}[label=(\arabic*)]
  \item 靠墙的一边不需篱笆，所以平行于墙的一边长为 $(60 - 2x)$ 米。\\
    面积 $S = x(60 - 2x) = -2x^2 + 60x$。\\
    定义域：$x > 0$ 且 $60 - 2x > 0$，即 $0 < x < 30$。
  \item $S = -2x^2 + 60x = -2(x^2 - 30x) = -2(x - 15)^2 + 450$。\\
    对称轴 $x = 15 \in (0, 30)$，$a = -2 < 0$，所以 $x = 15$ 时 $S$ 最大。\\
    $S_{\max} = 450$（平方米）。
\end{enumerate}
\textbf{答}：垂直于墙的一边长 $15$ 米时面积最大，为 $\boxed{450}$ 平方米。

\begin{appbox}[【应用3】基本不等式解决实际问题]
某工厂要建造一个容积为 $750$ 立方米的长方体无盖蓄水池。已知底面的造价为每平方米 $200$ 元，侧面的造价为每平方米 $150$ 元。设底面为边长 $x$ 米的正方形。
\begin{enumerate}[label=(\arabic*)]
  \item 写出总造价 $C$ 关于 $x$ 的函数。
  \item 当底面边长 $x$ 为多少时，总造价最低？最低造价是多少？
\end{enumerate}
\end{appbox}

\textbf{解}：
\begin{enumerate}[label=(\arabic*)]
  \item 底面面积 $= x^2$，水池高 $= \dfrac{750}{x^2}$。\\
    底面造价 $= 200x^2$。\\
    四个侧面面积 $= 4 \times x \times \dfrac{750}{x^2} = \dfrac{3000}{x}$。\\
    侧面造价 $= 150 \times \dfrac{3000}{x} = \dfrac{450000}{x}$。\\
    总造价 $C = 200x^2 + \dfrac{450000}{x}$（$x > 0$）。
  \item 用基本不等式。注意到 $C = 200x^2 + \dfrac{450000}{x}$，需要凑定值。\\
    将 $\dfrac{450000}{x}$ 拆成 $\dfrac{225000}{x} + \dfrac{225000}{x}$：\\
    $C = 200x^2 + \dfrac{225000}{x} + \dfrac{225000}{x}$。\\
    由三元均值不等式：$200x^2 + \dfrac{225000}{x} + \dfrac{225000}{x} \geq 3\sqrt[3]{200x^2 \times \dfrac{225000}{x} \times \dfrac{225000}{x}}$。\\
    $= 3\sqrt[3]{200 \times 225000 \times 225000} = 3\sqrt[3]{200 \times 225000^2}$。

    \textbf{换用配凑法}：令 $C = 200x^2 + \dfrac{450000}{x}$，设 $C' = 0$ 求极值点。\\
    $C'(x) = 400x - \dfrac{450000}{x^2} = 0 \Rightarrow x^3 = \dfrac{450000}{400} = 1125 \Rightarrow x = \sqrt[3]{1125} = \sqrt[3]{125 \times 9} = 5\sqrt[3]{9}$。

    用基本不等式验证：将 $200x^2$ 视为两项 $100x^2 + 100x^2$：\\
    $C = 100x^2 + 100x^2 + \dfrac{450000}{x} \geq 3\sqrt[3]{100x^2 \times 100x^2 \times \dfrac{450000}{x}} = 3\sqrt[3]{100^2 \times 450000 \times x^3}$。\\
    这个方法无法直接凑出常数，改用求导法。

    $C(x) = 200x^2 + \dfrac{450000}{x}$，$x > 0$。\\
    $C'(x) = 400x - \dfrac{450000}{x^2} = \dfrac{400x^3 - 450000}{x^2}$。\\
    $C'(x) = 0 \Rightarrow x^3 = 1125 \Rightarrow x = \sqrt[3]{1125} = 5\sqrt[3]{9} \approx 10.4$。\\
    $C''(x) = 400 + \dfrac{900000}{x^3} > 0$，确为最小值。\\
    $C_{\min} = 200 \times (5\sqrt[3]{9})^2 + \dfrac{450000}{5\sqrt[3]{9}} = 200 \times 25\sqrt[3]{81} + \dfrac{90000}{\sqrt[3]{9}} = 5000\sqrt[3]{81} + 90000 \times \dfrac{1}{\sqrt[3]{9}}$。\\
    化简：$5000 \times 9^{4/3}/9^{1/3} = 5000 \times 9 = 45000$... 让我们直接数值计算：\\
    $x = 5\sqrt[3]{9} \approx 10.40$，$C = 200 \times 108.16 + \dfrac{450000}{10.40} \approx 21632 + 43269 = 64901$。

    最低造价约为 $\boxed{64901}$ 元（当底面边长约为 $10.4$ 米时）。
\end{enumerate}

% =====================================================
% 第四部分：课堂练习
% =====================================================
\section*{\textcolor{tipgreen}{四、课堂练习}}

\begin{enumerate}[label=\textbf{\arabic*.}, leftmargin=2em]
  \item 解不等式 $x^2 - 4x + 3 > 0$。

  \item 解不等式 $-x^2 + 6x - 9 \geq 0$。

  \item 解不等式 $2x^2 + x - 1 \leq 0$。

  \item 已知 $x > 0$，求 $2x + \dfrac{8}{x}$ 的最小值。

  \item 已知 $x > 0,\ y > 0$，且 $2x + 3y = 12$，求 $xy$ 的最大值。

  \item 若不等式 $x^2 + mx + 1 > 0$ 的解集为 $\mathbb{R}$，求 $m$ 的取值范围。

  \item 已知方程 $2x^2 - 3x - 1 = 0$ 的两根为 $x_1, x_2$，不解方程，求 $x_1^2 + x_2^2$ 的值。

  \item 某矩形的周长为 $20$ 厘米，设一边长为 $x$ 厘米。
    \begin{enumerate}[label=(\alph*)]
      \item 写出面积 $S$ 关于 $x$ 的函数。
      \item 求面积的最大值。
    \end{enumerate}
\end{enumerate}

\vspace{1em}
\begin{knowledgebox}[参考答案]
\begin{enumerate}[label=\textbf{\arabic*.}, leftmargin=2em]
  \item $x^2 - 4x + 3 = (x-1)(x-3) > 0$，$a = 1 > 0$，$\Delta > 0$。解集为 $(-\infty,\ 1) \cup (3,\ +\infty)$。

  \item 乘 $-1$：$x^2 - 6x + 9 \leq 0$，即 $(x-3)^2 \leq 0$。只有 $x = 3$ 时等号成立。解集为 $\{3\}$。

  \item $\Delta = 1 + 8 = 9$，$x = \dfrac{-1 \pm 3}{4}$，$x_1 = -1$，$x_2 = \dfrac{1}{2}$。$a = 2 > 0$，"$\leq 0$"取中间。解集为 $\left[-1,\ \dfrac{1}{2}\right]$。

  \item $2x + \dfrac{8}{x} \geq 2\sqrt{2x \cdot \dfrac{8}{x}} = 2\sqrt{16} = 8$。等号在 $2x = \dfrac{8}{x}$ 即 $x = 2$ 时成立。最小值为 $8$。

  \item $2x + 3y = 12$，由基本不等式 $12 = 2x + 3y \geq 2\sqrt{6xy}$，所以 $xy \leq \dfrac{144}{24} = 6$。等号在 $2x = 3y = 6$ 即 $x = 3,\ y = 2$ 时成立。最大值为 $6$。

  \item 需要 $a = 1 > 0$ 且 $\Delta = m^2 - 4 < 0$，即 $-2 < m < 2$。

  \item $x_1 + x_2 = \dfrac{3}{2}$，$x_1 x_2 = -\dfrac{1}{2}$。$x_1^2 + x_2^2 = (x_1 + x_2)^2 - 2x_1 x_2 = \dfrac{9}{4} + 1 = \dfrac{13}{4}$。

  \item \begin{enumerate}[label=(\alph*)]
    \item 另一边为 $(10 - x)$，$S = x(10 - x) = -x^2 + 10x$，$0 < x < 10$。
    \item $S = -(x-5)^2 + 25$，当 $x = 5$ 时 $S_{\max} = 25$ 平方厘米。
  \end{enumerate}
\end{enumerate}
\end{knowledgebox}

\vfill
\begin{center}
\small\textcolor{gray}{本讲义仅供教学参考，例题可根据学生水平调整难度。}
\end{center}

\end{document}
